% !TEX root = ../journal_v2_0421052099.tex
%
% Condensed from thesis_book/chapters/framework.tex. Checked against the code
% in agr/ (nodes.py, scorer.py, planner.py, budget.py, kg_tools.py).
\section{The AGR Framework}\label{sec:framework}

\subsection{The Controlled-Agent Pattern}

There are two ways to build a graph-navigating agent. The first hands the model
a set of tools through native function calling and lets it drive the loop
itself, in the style of ReAct~\cite{yao2023react}. The second inverts control.
The \emph{program} orchestrates the loop, calls the graph tools
deterministically, and consults the model only at decision points, through
constrained prompts that must return JSON. AGR uses the second. Every graph
operation is then logged by construction, and budgets are enforced in code
rather than requested in prompts. A malformed model response can spoil one
decision, but it cannot corrupt the control flow. The agency in AGR lives in the
decision policy. That policy decides which sub-objectives to pursue, which edges
look promising, whether a path is a dead end, and whether the draft answer is
supported.

\subsection{The State Machine}

AGR is a state machine of six nodes over a shared typed state, built with
LangGraph (\Cref{fig:statemachine}). The topology has cycles. Exploration loops
until the evaluator is satisfied, and verification can send the agent back to
explore. The verifier sits between the decision to answer and the answer itself.
There is no path from the evaluator to the answerer that does not pass through
verification.

\begin{figure}[!t]
\centering
% !TEX root = ../journal_v2_0421052099.tex
%
% Hand-drawn. The AGR state machine of the book's Figure 4.1
% (thesis_book/chapters/framework.tex), redrawn narrower so that it fits one
% journal column. Topology, labels and colours are the book's. Blue marks the
% agent nodes and green the verifier.
\begin{tikzpicture}[
    node distance=5mm,
    box/.style={draw=agrNode, thick, rounded corners=2pt, minimum height=6.5mm,
      fill=agrNode!8, minimum width=21mm, align=center, font=\footnotesize},
    vbox/.style={box, draw=agrSup, fill=agrSup!12},
    term/.style={draw=black!45, rounded corners=5pt, minimum height=5mm,
      minimum width=14mm, font=\scriptsize\sffamily, fill=black!6},
    flow/.style={-{Stealth[length=1.8mm]}, semithick, draw=black!68},
    lbl/.style={font=\scriptsize\ttfamily, inner sep=1.2pt, fill=white,
      anchor=center}]

  \node[term]              (start) {START};
  \node[box, below=of start] (plan) {Planner};
  \node[box, below=of plan]  (expl) {Explorer};
  \node[box, below=of expl]  (eval) {Evaluator};
  \node[vbox, below=of eval] (ver)  {Verifier};
  \node[box, below=of ver]   (ans)  {Answerer};
  \node[term, below=of ans]  (end)  {END};
  \node[box, right=17mm of eval, minimum width=19mm] (back) {Backtracker};

  \draw[flow] (start) -- (plan);
  \draw[flow] (plan) -- (expl);
  \draw[flow] (expl) -- (eval);
  \draw[flow] (eval) -- node[lbl] {answer} (ver);
  \draw[flow] (ver) -- node[lbl] {grounded $\mid$ give\_up} (ans);
  \draw[flow] (ans) -- (end);

  % The search cycle, from the evaluator back to the explorer.
  \draw[flow] (eval.west) -- ++(-15mm,0) node[lbl, above, pos=0.5]
    {continue} |- (expl.west);

  % The backtracking cycle, through the backtracker.
  \draw[flow] (eval.east) -- node[lbl] {backtrack} (back.west);
  \draw[flow] (back.north) |- (expl.east);

  % Verification-driven re-exploration.
  \draw[flow] (ver.west) -- ++(-21mm,0) node[lbl, below, pos=0.5] {retry} |-
    ([yshift=-1.8mm] expl.west);

  \begin{scope}[on background layer]
    \node[draw=black!35, dashed, rounded corners, inner sep=2mm,
          fit=(expl)(eval)] {};
  \end{scope}
\end{tikzpicture}

\caption{The AGR state machine. The dashed region is the explorer--evaluator
  search loop. The backtracker restores an earlier frontier, and the verifier's
  retry edge is verification-driven re-exploration. All three cycles are bounded
  by the budgets of \Cref{tab:budget}. Every path from the evaluator to the
  answerer passes through the verifier.}
\label{fig:statemachine}
\end{figure}

\subsection{The Graph Tool API}\label{sec:tools}

The model never writes graph queries. A generated query introduces a failure
that is indistinguishable, in the results, from a reasoning failure. When a
question is answered wrongly, one cannot tell whether the agent reasoned badly
or emitted invalid syntax. So the tools are parameterised Cypher templates,
written once and unit-tested. Four operations are live.

\textsf{search\_entity} links a mention to candidate nodes through the resolver
of \Cref{sec:environment}. It is the only entry point from text to graph.

\textsf{get\_relations} returns the relation types incident to an entity in
\emph{both} directions, each with its fanout count. Freebase's edge direction is
a modelling artefact. And the count is free signal, since a relation with a
fanout of forty thousand rarely answers a specific question. Relations in
administrative namespaces such as \textsf{common}, \textsf{type}, and
\textsf{freebase} are dropped. The list is capped at $300$ relations, sorted by
fanout.

\textsf{get\_neighbors} expands an entity along one relation. When it reaches a
mediator, it hops through it and returns the entities on the far side, each with
the chain of relations that reached it. So the model reasons about logical
neighbours, while the log keeps the raw chain. At most $200$ neighbours are
returned per expansion.

\textsf{verify\_connection} asks whether two entities are adjacent in the full
graph, directly or through one mediator. It ignores the relation and the
direction. A fifth operation, \textsf{verify\_triple}, also matched the
relation. It was built first as the claim checker and did not survive contact
with generated text. Claim decomposition produces relations phrased in English,
while Freebase predicates are neither phrased that way nor reliably one edge
long. So requiring the relation to match rejected claims that were true, and no
node calls it.

\subsection{Planner, Explorer, and Scorer}\label{sec:explorer}

The planner makes one model call. It turns the question into an ordered chain of
at most four sub-objectives, in which later objectives refer back to earlier
results with a \textsf{\#N} back-reference. The prompt is few-shot, with one
example for each of four archetypes (single hop, composition, conjunction, and
superlative). The single-hop example shows that \emph{not} decomposing is a
valid output. The plan is validated rather than trusted. If it is empty, too
long, refers forwards, or names no topic, the planner falls back to a single
sub-objective holding the whole question. Planning can degrade but can never
abort a run. Search is seeded from the dataset's annotated topic entities, the
standard protocol in this literature~\cite{sun2024think,luo2024reasoning}. That
choice separates reasoning from entity linking.

The explorer performs one expansion of the frontier. It pushes a checkpoint onto
the backtrack stack and gathers the relations of every anchor into a single
candidate list. It scores the whole list in one pass and takes the best three
edges per anchor, skipping banned ones. Each selected edge is expanded, and the
new frontier is sorted by the score of the edge that produced each entity and
truncated to five anchors. Each candidate edge $c$ receives the hybrid score
\begin{equation}\label{eq:score}
\mathrm{score}(c) = \alpha \cos\big(\mathbf{v}_{\mathrm{rel}(c)},
\mathbf{v}_{\mathrm{obj}}\big) + (1-\alpha)\,\mathrm{LLM}(c),
\end{equation}
where $\mathbf{v}_{\mathrm{rel}(c)}$ embeds the candidate's relation and
$\mathbf{v}_{\mathrm{obj}}$ embeds the active sub-objective. $\mathrm{LLM}(c)
\in [0,1]$ is the model's rating of how likely following the edge is to advance
that objective. Three properties matter. The objective is embedded, not the
question, with any back-reference already replaced by the resolved entity's
name. That substitution is the mechanism by which decomposition is supposed to
help. The model rates only the twenty candidates that the embedding ranks
highest, in one batched call. Every other candidate scores $\mathrm{LLM}(c) =
0$, so the pre-filter is a hard gate. And the blend is applied after the call.
The scoring prompt holds no configuration value, so every condition of an
$\alpha$ sweep shares the same cached model responses. We use $\alpha = 0.7$,
chosen on development data (\Cref{sec:ablation-setup}).

\subsection{Evaluator and Backtracking}

The evaluator makes one model call per expansion. It sees the active
sub-objective, the thirty traversed facts most relevant to it, and the remaining
budget. It decides whether to continue, backtrack, or answer, and it names the
entities that satisfy the objective. Those names are not taken on trust. Any
name that is not the tail of a traversed triple is discarded. A language model
asked which entities satisfy an objective will happily list ones it knows from
training but never retrieved. This \emph{groundedness filter} refuses to let
that recall pass as evidence. On one question of an early smoke run, the
evaluator listed nine correct countries that appeared in no retrieved triple,
because the beam had gone down an unrelated branch. The filter rejected all
nine, and the system abstained. That cost a point, but it kept every proposed
entity tied to the graph. Surviving resolutions are added to the candidate
answers on every evaluation, whether or not the objective was judged complete.
When an objective completes, the plan advances and the resolved entities become
the next frontier.

Backtracking fires on the first of three triggers, checked in this order. Either
the last expansion produced no new triples, or the frontier's signal fell below
a threshold $\tau = 0.20$, or the evaluator judged the direction a dead end. The
first two cost nothing, so a mechanically hopeless frontier is abandoned without
asking the model. The signal compared against $\tau$ is the largest of three
quantities. They are the best blended score among the expanded edges, the best
raw embedding similarity over every candidate, and the best model rating. So the
trigger fires only when all three are weak. The backtracker restores the
\emph{highest-scoring} snapshot on the stack rather than the most recent one. It
also bans the (anchor, relation, direction) combinations of the last pass.
Without that ban list, a deterministic scorer facing a restored frontier would
select the same edges again until the budget ran out. \Cref{alg:nav} states the
whole navigation loop.

\begin{algorithm}[!t]
\caption{AGR navigation, with the verification layer abstracted}
\label{alg:nav}
\small
\begin{algorithmic}[1]
\REQUIRE question $q$, budgets, blend $\alpha$, threshold $\tau$
\ENSURE answer $A$ with supporting triples $S$
\STATE $P \gets \textsc{Plan}(q)$; \textbf{if} invalid \textbf{then}
  $P \gets [\,q\,]$
\STATE $F \gets \textsc{SeedAnchors}(q)$; $T, C, \mathit{stack},
  \mathit{ban} \gets \emptyset$
\WHILE{budgets permit}
  \STATE push $(F, \mathit{depth}, \sigma)$ onto $\mathit{stack}$
  \STATE $R \gets \bigcup_{e \in F} \textsc{GetRelations}(e)$
  \STATE score $R$ by \Cref{eq:score} against the active objective
  \STATE $E \gets$ best $3$ edges per anchor, excluding $\mathit{ban}$
  \STATE $T \gets T \cup \textsc{GetNeighbors}(E)$
  \STATE $F \gets$ best $5$ new entities by edge score
  \STATE $\sigma \gets$ the largest of the three signals of this pass
  \STATE $(d, \mathit{done}, \rho) \gets \textsc{Evaluate}(\text{objective},
    \textsc{TopFacts}(T))$
  \STATE $\rho \gets \rho \cap \textsc{Tails}(T)$; $C \gets C \cup \rho$
    \COMMENT{groundedness filter}
  \IF{$\mathit{done}$}
    \STATE advance $P$; \textbf{if} no objective remains \textbf{then}
      $d \gets \textsf{answer}$ \textbf{else} $F \gets \rho$
  \ENDIF
  \IF{no new triples \OR $\sigma < \tau$ \OR $d = \textsf{backtrack}$}
    \STATE $\mathit{ban} \gets \mathit{ban} \cup E$; restore the
      highest-scoring snapshot
  \ELSIF{$d = \textsf{answer}$}
    \STATE \textbf{break}
  \ENDIF
\ENDWHILE
\STATE $(A, S) \gets \textsc{VerifyAndAnswer}(q, T, C)$
  \COMMENT{\Cref{alg:verify}}
\RETURN $(A, S)$
\end{algorithmic}
\end{algorithm}

The implementation departs from this specification in one way that has a
measurable consequence. It is recorded rather than repaired, since repairing it
would change the frozen results. The backtracker bans the \emph{most recent}
pass but restores the \emph{highest-scoring} snapshot, which is often older. So
whatever was expanded in between stays unbanned, and the explorer can select it
again. Over the main test runs, $53$ explorer passes expanded a set of (anchor,
relation) pairs that the same question had already expanded. That is exactly the
repetition the ban list exists to prevent.

\subsection{Budgets and Termination}

\Cref{tab:budget} lists the budgets. They are enforced in the routers between
nodes and in the model client, never in prompts. The depth budget counts
explorer expansions. One expansion can cross a mediator, so it does not bound
the search at four raw hops. Every cycle in \Cref{fig:statemachine} passes
through a router that checks a counter which only moves one way. So no sequence
of model outputs can produce a run that never ends. Every model call is also
cached, keyed on the model, its sampling settings, and the full prompt. A cache
hit replays the original token counts, so a cached rerun reproduces the original
run's budget record exactly. Each run writes two logs. The tool log records
every graph operation with its arguments, result size, and latency. The run
record holds the answer, the draft, the verification outcome, every backtrack
with its trigger, the full decision trace, and the budget snapshot. It also
names the backbone, the budget configuration hash, the run configuration, and
the source commit. So any record found anywhere states what produced it.

\begin{table}[!t]
\caption{Budget configuration. The last two rows are caps in the tool layer that
  bound what the agent can see.}
\label{tab:budget}
\centering
\footnotesize
\begin{tabular}{@{}lrl@{}}
\toprule
Budget & Value & What it bounds \\
\midrule
\textsf{max\_depth} & $4$ & explorer expansions \\
\textsf{beam\_width} & $3$ & edges expanded per anchor \\
\textsf{max\_anchors} & $5$ & frontier size \\
\textsf{max\_backtracks} & $3$ & backtracks per question \\
\textsf{max\_llm\_calls} & $25$ & model calls per question \\
\textsf{max\_verify\_iters} & $2$ & repair iterations \\
\textsf{max\_seconds} & $300$ & wall-clock time \\
\midrule
\textsf{max\_relations} & $300$ & relations offered to the scorer \\
\textsf{max\_fanout} & $200$ & neighbours per expansion \\
\bottomrule
\end{tabular}
\end{table}

\subsection{Design Validation on a Smoke Set}\label{sec:smoke}

Twenty questions from the training and validation splits, never from test, were
used once to check the assembled framework's mechanics. The set was built to
exercise specific machinery rather than to sample the benchmarks. The first
end-to-end run answered roughly three or four of them correctly, and the final
configuration answers thirteen or fourteen. That interval is where several of
the mechanisms above were shown to be necessary. The first run's failures
collapsed into five systematic causes, and \Cref{tab:rootcauses} gives the
mechanism each one motivated. Its before-state was read from a console log that
was never committed. So these figures are given as ranges, and none of them
enters any result.

\begin{table}[!t]
\caption{Root causes of failure in the first end-to-end run on the smoke set,
  and the mechanism each one motivated.}
\label{tab:rootcauses}
\centering
\footnotesize
\begin{tabular}{@{}L{43mm}L{39mm}@{}}
\toprule
Observed failure & Mechanism introduced \\
\midrule
Backtracking was a no-op, since a restored frontier reproduced the same scores and decisions & A ban list on abandoned (anchor, relation, direction) combinations \\
\addlinespace
The fact window showed the thirty most recent triples, so a large expansion pushed the answer out of view & A fact window ranked by relevance to the objective \\
\addlinespace
A relation cap of $80$, sorted by fanout, hid precise low-fanout relations behind hub relations & The cap raised to $300$, and administrative predicates blocked \\
\addlinespace
Entities resolved while an objective was incomplete were discarded & Candidates accumulated on every evaluation \\
\addlinespace
The frontier was truncated in insertion order, so the agent drifted from the anchor that mattered & The frontier sorted by edge score before truncation \\
\bottomrule
\end{tabular}
\end{table}

The two confabulations of the first run came from a retrieval defect, not a
generation defect. Two questions about national currencies produced a confident
``United States Dollar''. The relation cap had hidden the currency relation, the
beam surfaced triples about national income per capita in dollars, and the
answerer matched the pattern. No retrieved triple supported the answer, and
nothing that examined the answer caught it. After the cap was raised, both
questions resolved correctly at roughly a third of the token cost. But a
claim-level check against the traversed triples would have caught them
regardless of the retrieval defect. That is the argument for verifying before
emitting.

\subsection{The Structural Verification Layer}\label{sec:verification}

The navigation loop ends with a set of traversed triples and a list of candidate
entities. Nothing in it constrains what the final answer \emph{asserts}. A final
generation call is free to name an entity that appeared in neither, as the
currency answers of \Cref{sec:smoke} did.

The response is to check the draft before it is emitted. Chain-of-Verification
and FActScore decompose and check after generation, and both take further model
output as evidence~\cite{dhuliawala2024chain,min2023factscore}. The layer
borrows their decomposition step and rejects their placement and their evidence.
Placing the check inside the loop has three consequences. The checker sees the
same evidence the drafter saw. A rejected claim can trigger more exploration
while budget remains. And the judge of whether a fact holds is the environment,
not a second opinion from the model that wrote the draft.

\textbf{Claim decomposition.} One model call produces three things at once. They
are a prose draft, the entities the draft asserts as the answer, and a
decomposition of the draft into atomic claims $\langle h, r, t \rangle$ over
entity names. The call sees at most sixty traversed triples, ranked by relevance
to the question, and up to twenty-five candidates. The prompt requires names to
be copied exactly from the facts, and it requires a hedged draft to carry no
claims. It also requires the answer entities to be specific entities of the kind
the question asks for. That rule was added after drafts such as ``the senators
from New Jersey are the United States Senate members associated with New
Jersey''. Such an answer is grammatically complete and semantically empty. The
relation slot $r$ is free text, since the model has no reliable way to name a
Freebase predicate. Decomposition is sparse. On the development set the drafter
emitted $121$ claims over $80$ questions, with a median of one per question. Ten
questions produced no claims at all, so for those ten the layer certified
nothing.

\begin{figure*}[!t]
\centering
% !TEX root = ../journal_0421052099.tex
%
% Copied unchanged from thesis_book/figures/fig_claim_path.tex. Hand-drawn (not
% generated): the path one atomic claim takes through the verification layer.
% Mirrors Algorithm 1 line for line, including which of the three routes
% records citable evidence.
%
% Colour is semantic, not decorative: blue marks a step that does work, green
% the supported bucket, vermillion the unsupported one. The palette is the same
% colour-blind-safe set the generated figures use, and every distinction is
% still carried by shape and position, so a greyscale print loses nothing.
% Defined here because this figure is hand-drawn and has no generator.
\definecolor{agrNode}{HTML}{0072B2} \definecolor{agrSup}{HTML}{009E73}
\definecolor{agrUns}{HTML}{D55E00}
% This figure is set at two different sizes in two different text fonts (the
% thesis body font and the slide font), and the node widths are in millimetres
% either way. Hyphenation is switched off inside the picture so a node cannot
% break a word in one document and not the other; widths below are the ones that
% hold every label on the line count it was drawn for in both.
\begin{tikzpicture}[
    node distance=6mm,
    execute at begin picture={\hyphenpenalty=10000\exhyphenpenalty=10000},
    box/.style={draw=agrNode, thick, rounded corners=2pt, minimum height=7mm,
                text width=33mm, align=center, font=\small,
                fill=agrNode!8},
    dec/.style={draw=agrNode!55, rounded corners=2pt, minimum height=8mm,
                text width=38mm, align=center, font=\scriptsize,
                fill=agrNode!4},
    bucket/.style={draw, rounded corners=6pt, minimum height=6mm,
                text width=21mm, align=center, font=\scriptsize\sffamily},
    sup/.style={bucket, draw=agrSup, thick, fill=agrSup!14},
    uns/.style={bucket, draw=agrUns, thick, fill=agrUns!14},
    flow/.style={-{Stealth[length=2mm]}, semithick, draw=black!68},
    lbl/.style={font=\scriptsize, inner sep=1.5pt, fill=white},
    ev/.style={font=\scriptsize\itshape, inner sep=1.5pt, text=agrSup}]

  \node[box]                  (draft) {Draft and decompose
                                       \\[-2pt]{\scriptsize one model call}};
  \node[dec, below=of draft]  (d1)    {both endpoint names seen
                                       during traversal?};
  \node[dec, below=of d1]     (d2)    {pair in \emph{traversed}
                                       adjacency?};
  \node[dec, below=of d2]     (d3)    {\textsf{verify\_connection}:
                                       adjacent in the full graph?};
  \node[box, below=9mm of d3] (ent)   {entailment check
                                       \\[-2pt]{\scriptsize one batched call}};

  \node[sup, right=30mm of d2] (sup)  {supported $S$};
  \node[uns, right=30mm of ent] (uns) {unsupported $U$};

  \draw[flow] (draft) -- node[lbl, pos=0.5] {claims $K$} (d1);
  \draw[flow] (d1) -- node[lbl] {yes} (d2);
  \draw[flow] (d2) -- node[lbl] {no} (d3);

  % The one route that records evidence. "+ evidence" is the widest label on
  % the picture; at pos=0.80 its right edge crossed into the supported bucket,
  % so it is pulled back to the middle of the run instead.
  \draw[flow] (d2.east) -- node[lbl, pos=0.24] {yes}
              node[ev, pos=0.66, above] {+ evidence} (sup.west);
  % Arrives from BELOW, not along the bucket's own bottom edge. As "|- to
  % sup.south east" this turned at exactly the height the entailed route
  % turns at, ran on top of it as far as the bucket, then continued
  % underneath the bucket to land on its far corner: two arrows drawn over
  % each other and one of them apparently passing through a box.
  \draw[flow] (d3.east) -- node[lbl, pos=0.28] {yes} ++(16mm,0)
              -| ([xshift=8mm]sup.south west);
  \draw[flow] (d3) -- node[lbl] {no} (ent);
  % Endpoint never traversed: skip both structural routes entirely. The run is
  % long enough to clear the "structural routes" caption, which sits to the
  % left of the dashed region at the same height as this line's descent.
  \draw[flow] (d1.west) -- ++(-26mm,0) node[lbl, pos=0.5] {no} |- (ent.west);

  \draw[flow] (ent.east) -- ++(9mm,0) coordinate (e)
              |- node[lbl, pos=0.75] {entailed} (sup.south west);
  \draw[flow] (e) -- (uns.west);

  \node[bucket, draw=black!45, fill=black!5, below=11mm of uns.south,
        anchor=north, text width=56mm]
        (route) {$U = \emptyset \Rightarrow$ \textsf{grounded}\quad$\mid$\quad
                 budget $\Rightarrow$ \textsf{retry}\quad$\mid$\quad
                 \textsf{give\_up}};
  \draw[flow] (uns) -- (route);
  % Down into the routing box's top edge, near its right corner.
  % "-- ++(7mm,0) |- (route.east)" put the descent 7mm right of the bucket,
  % which is still well inside the routing box: the line came down through
  % the box's interior and the arrowhead sat on its border pointing
  % outwards, which is what it looked like -- an arrow leaving the box.
  % Turning down INSIDE route.east rather than outside it keeps the fix
  % from widening the picture; routed round the outside instead, the paper
  % -- which sets this figure unscaled -- went 25.8pt overfull.
  \draw[flow] (sup.east) -| ([xshift=-6mm]route.north east);

  \begin{scope}[on background layer]
    \node[draw=black!35, dashed, rounded corners, inner sep=2.5mm,
          fit=(d2)(d3), label={[font=\scriptsize\itshape, anchor=west,
          xshift=1mm]west:}] (strbox) {};
  \end{scope}
  \node[font=\scriptsize\itshape, left=1.5mm of strbox, align=right,
        text width=17mm] {structural\\routes};
\end{tikzpicture}

\caption{The path of one atomic claim through the verification layer, mirroring
  \Cref{alg:verify}. A claim leaves at the first test that settles it. Only the
  first structural route adds to the evidence, which is why some grounded
  answers carry supporting triples and others do not. A claim that names an
  entity the agent never traversed skips both structural routes, because the
  name map is built from traversed triples alone.}
\label{fig:claim}
\end{figure*}

\textbf{Two-stage checking.} \Cref{alg:verify} gives the procedure, and
\Cref{fig:claim} draws one claim's path through it. The structural check first
maps a claim's names back to nodes through the traversed triples alone. A claim
naming an entity the agent never traversed cannot reach either structural route.
It falls to the entailment check. The first structural route is \emph{traversed
adjacency}. The claim is supported if a traversed triple joins its endpoints,
and every such triple is attached to the answer as evidence. The second route
asks \textsf{verify\_connection} whether the endpoints are adjacent in the full
graph, since the agent may have reached both along different branches. It
records no evidence. This second route is rare. On the test sets it was
consulted for $100$ of the $2{,}110$ claims and accepted $39$. Claims that fail
both routes go to the \emph{entailment check}, one batched model call that
judges each claim against the same fact window. If that call fails for any
reason, every pending claim is marked unsupported.

The structural check runs first for three reasons. The graph should decide where
it can, so the model is consulted only where the graph is silent. Only traversed
adjacency yields triples that can be attached to the answer. And a set lookup is
free, while the entailment call is the layer's only unavoidable cost beyond
drafting.

\begin{algorithm}[!t]
\caption{Structural verification of a draft}
\label{alg:verify}
\small
\begin{algorithmic}[1]
\REQUIRE question $q$, traversed triples $T$, candidates $C$
\ENSURE draft, answer entities $A$, claim sets $S$ and $U$, evidence $E$
\STATE $W \gets \textsc{TopFacts}(q, T, 60)$
\STATE $\mu \gets$ name-to-id map over $T$
\STATE $(\mathit{draft}, A, K) \gets \textsc{DraftAndDecompose}(q, W, C)$
\STATE $\mathit{adj} \gets \{(h,t),(t,h) : (h,r,t) \in T\}$
\STATE $S \gets \emptyset$; $P \gets \emptyset$; $E \gets \langle\,\rangle$
\FORALL{$c = \langle h, r, t \rangle \in K$}
  \IF{$h \notin \mu$ \OR $t \notin \mu$}
    \STATE $P \gets P \cup \{c\}$ \COMMENT{never traversed}
  \ELSIF{$(\mu[h], \mu[t]) \in \mathit{adj}$}
    \STATE $S \gets S \cup \{c\}$; append the joining triples to $E$
  \ELSIF{$\textsf{verify\_connection}(\mu[h], \mu[t])$}
    \STATE $S \gets S \cup \{c\}$ \COMMENT{no evidence recorded}
  \ELSE
    \STATE $P \gets P \cup \{c\}$
  \ENDIF
\ENDFOR
\STATE $U \gets \emptyset$
\IF{$P \neq \emptyset$}
  \STATE split $P$ into $S$ and $U$ by $\textsc{Entail}(W, P)$
\ENDIF
\IF{$U \neq \emptyset$ \AND repair budget remains}
  \STATE add a repair objective for the first claim in $U$
\ENDIF
\STATE route \textsf{grounded} if $U = \emptyset$, else \textsf{retry}
  if the repair, depth, and time budgets remain, else \textsf{give\_up}
\RETURN $(\mathit{draft}, A, S, U, E)$
\end{algorithmic}
\end{algorithm}

\textbf{Repair.} The router after the verifier has three outcomes. With no
unsupported claim it routes to \textsf{grounded}, and the draft is emitted with
no further model call. With unsupported claims and the repair, depth, and time
budgets all remaining, it routes to \textsf{retry}. The verifier has already
appended a repair objective for the first unsupported claim and re-anchored the
search on that claim's head entity. Anything else routes to \textsf{give\_up}.
On that path the final entity list is computed deterministically. An entity is
kept if it appears in a supported claim or in no unsupported claim. One rewrite
call then restates the prose around the kept list. Repair is capped at two
iterations. The grounded path costs nothing, so the layer is cheap overall. On
the development set it raised mean tokens per question from $4{,}858$ to
$4{,}925$, about $1.4\%$.

The retry route never ran on the development set. On all three development
questions that produced an unsupported claim, the depth budget was already spent
when the verifier was reached. That correlation does not survive the move to
test. Across the $800$ test questions the verifier asked for a repair on $52$,
and exploration actually resumed on $28$ of them. So the depth condition blocks
the retry about half the time. Thirteen questions ended grounded after a repair,
and the iteration cap bound on $15$.

That deterministic filter replaced a second generation call, and the reason is
the most instructive result of the design. On the development set, a
question-by-question comparison against the same system with claim checking
switched off found exactly three questions that differed. No entailment verdict
was involved in any of them. All three traced to the rewrite call, which at the
time regenerated the answer entities along with the prose. In one case it
dropped four correct entities together with an unsupported one. In another it
added the question's own subject to the answer. In the third it respelled the
graph's \textsf{University Yale} as \textsf{Yale University}, which scored as
wrong (\Cref{sec:example}). The verification logic was sound, and its repair
path frequently was not. Replacing generation with filtering removed two of the
three regressions.

\textbf{The output contract and its limits.} AGR returns each answer paired with
the triples that support it. That coverage is narrower than the phrase suggests.
Only the traversed-adjacency route records evidence. And the run log keeps the
number of supporting triples, not the triples. So the contract holds inside the
run, but it cannot be checked afterwards from the committed records. On the
development set, $13$ of $80$ answers carry no supporting triples at all. Ten of
them are hedges with no claims, for which zero is the honest count. The other
three had every claim rejected or certified only through the route that records
nothing.

\subsection{A Worked Example}\label{sec:example}

The question ``Which university, founded in 1701, did actor James Franco
attend?'' shows the layer end to end. Its gold answer is the graph's
\textsf{University Yale}. Navigation anchors on James Franco and expands his
education relations. The plan's first objective, the university founded in 1701,
is never satisfied, because the environment holds essentially no date literals.
After three evaluator-triggered backtracks, the run stops on the depth budget.

The drafter writes ``James Franco attended University Yale, which was founded in
1701.'' Its entity list is \textsf{[University Yale]}, copied from the fact
window in the graph's spelling. It decomposes the draft into an attendance claim
and a founding claim. Traversed triples join the attendance claim's endpoints,
so that claim is supported, and four triples are recorded as its evidence. The
founding claim names \textsf{1701}, which no traversed triple mentions. So it
falls to the entailment check, which correctly reports it unsupported. The depth
budget is spent, so the router takes \textsf{give\_up}. The filter keeps
\textsf{University Yale}, since it appears in a supported claim. The rewrite
asserts the attendance and hedges the founding date, and the answer scores as a
hit.

The instructive part is the counterfactual. Before deterministic filtering, the
identical trace ended with \textsf{[Yale University]}. The rewrite call had
restated the answer in the model's own spelling rather than the graph's, and the
answer was scored wrong. The claim-level checks had worked perfectly. The layer
still emitted an entity string that appears nowhere in the environment, which is
the precise thing it exists to prevent.

\subsection{What the Layer Can Get Wrong}\label{sec:failure-modes}

The ways the layer can be wrong follow from its design, independently of how
often they occur. A claim is \emph{wrongly rejected} when the layer withholds
support that the environment would in fact provide. It is \emph{wrongly
accepted} when the layer certifies a claim that the environment does not
support. \Cref{sec:errors} reports which of these mechanisms actually fired.

Five mechanisms lead to wrongful rejection. The most expensive is the composite
claim. When a draft asserts a filtered set, such as ``these teams were founded
before 1996'', its claims are about the founding. The membership that makes each
team a candidate may never be claimed at all. So each team's only claim is the
unverifiable one, the filter drops every team, and a draft with four correct
answers becomes a hedge. A claim can also name an endpoint outside the
traversal, which skips both structural routes whatever the graph contains. The
entailment check's fact window is ranked for relevance to the question rather
than to the claim, so a peripheral claim may be judged against facts that never
mention it. A failed entailment call records an infrastructure error as a
rejection. And the name map keeps the first identifier it sees for a name, so
two entities that share a surface form collapse into one.

Wrongful acceptance has more mechanisms, and they reach further. Both structural
routes ignore the relation and the direction. So any edge between two entities
certifies any claimed relation between them. A claim that X is Y's mother is
accepted on an edge recording that X is Y's child. That is the price of
free-text relations, and it is the layer's main acceptance risk. On the test
sets the verifier accepted $2{,}008$ of the $2{,}110$ claims it decomposed. The
number accepted by a relation-blind test lies somewhere in $[39, 2{,}008]$. The
lower end is what the log pins down, and the upper end is what the mechanism
permits. The mediator hop adds a second looseness. An actor, a film, a
character, and a period are all mutually connected through one performance node,
so any claim relating two of them is certified. A draft with no claims has no
unsupported claim and is routed \textsf{grounded}, which describes $10$ of the
$80$ development questions. An entity that the drafter named but the decomposer
never put into a claim is never checked. So the guarantee over emitted entities
rests on decomposition recall, which is itself an unchecked model output. A
claim accepted structurally is never examined for meaning, since the entailment
check sees only what the structural routes could not resolve. And an answer
whose every claim is supported can still fail to answer the question, as the New
Jersey senators draft does. No claim-level check can catch that, because the
defect lies between the claims and the question.

The layer establishes that the entity pairs a draft relates are connected in the
environment. On the repair path it also makes the emitted entity list a
deterministic function of those verdicts rather than a second generation. It
does not establish that the claimed relation is the one the edge records, that
an entity absent from every claim was grounded, or that the answer is responsive
to the question. Those boundaries are properties of the design rather than
defects discovered in it.

\endinput
